% Transcription — fonds Alexandre Grothendieck, shelfmark 66, batch 1 (pages 1-4).
% Established from the facsimile archives/batches/66/batch-01.pdf (a mirror of
% https://grothendieck.umontpellier.fr/66.pdf), per the skill
% transcribe-grothendieck.
% A list of twenty-eight thesis subjects, dated 1964 in the author's own hand —
% one of the rare folders that carries its date rather than receiving one.
% Circled numbers are the ones he judged hard; he says so at the head of the
% page. The left margins carry later annotations in a different ink recording
% what became of a subject — who took it, what was abandoned, and in two cases
% the answer that arrived (Raynaud on 17, Mumford on 25).
% Pass: Opus 5 (claude-opus-5), 2026-08-16 — first pass, unchecked against the pages by a human.
\documentclass[11pt,a4paper]{article}
% Le préambule commun à toutes les transcriptions du fonds Grothendieck.
%
% Il tient en une idée : **l'incertitude se compose, elle ne se résout pas.**
% Une transcription qui lisse un mot douteux a détruit la seule chose pour
% laquelle on la faisait — un lecteur doit pouvoir savoir, à chaque endroit,
% ce qui était sur le feuillet et ce qui a été ajouté. Les six macros qui
% suivent sont donc l'appareil critique, et non de la décoration.
%
% Les mêmes commandes sont reconnues par `scripts/render.mjs`, qui produit la
% vue de lecture du site. Une macro ajoutée ici doit l'être aussi là-bas,
% faute de quoi le rendu échouera bruyamment — ce qui est voulu : un rendu
% silencieusement faux est pire qu'un rendu absent.

\ProvidesPackage{grothendieck}[2026/08/08 Transcriptions du fonds Grothendieck]

\usepackage{amsmath,amssymb,amsthm}
\usepackage{mathrsfs}
\usepackage{stmaryrd}
% Les diagrammes commutatifs sont partout dans ce fonds : c'est la langue
% même de la géométrie algébrique de Grothendieck, et une transcription qui
% les rendrait en prose ne transcrirait rien.
\usepackage{tikz-cd}
% Français par défaut — c'est la langue des feuillets — mais l'anglais est
% chargé aussi : la traduction et le résumé sont des documents anglais, et la
% typographie française (espaces avant les deux-points) y serait une faute.
% Ces documents basculent par \selectlanguage{english} en tête de corps.
\usepackage[english,french]{babel}
\usepackage{fontspec}
\usepackage{geometry}
\geometry{a4paper,margin=25mm,marginparwidth=28mm}
\usepackage{marginnote}
\usepackage{xcolor}
% ulem, not soul. `soul`'s \st and \ul are built on a character-by-character
% scan that XeLaTeX's font machinery breaks: the document fails with an
% undefined control sequence rather than degrading. ulem does the same two jobs
% and is Unicode-engine safe. `normalem` stops it hijacking \emph, which the
% transcriptions use for Grothendieck's own emphasis.
\usepackage[normalem]{ulem}
\usepackage{hyperref}
\hypersetup{colorlinks=true,linkcolor=black,urlcolor=black,pdfborder={0 0 0}}

% --- Métadonnées du lot -----------------------------------------------------
% Reprises telles quelles par le script de rendu, qui les affiche en tête de
% la vue de lecture. Elles fixent ce que le fichier prétend couvrir : sans
% elles, une transcription flotte sans point d'attache dans un fonds de seize
% mille feuillets.

\newcommand{\folder}[1]{\gdef\grothFolder{#1}}
\newcommand{\batch}[1]{\gdef\grothBatch{#1}}
\newcommand{\leaves}[2]{\gdef\grothFirst{#1}\gdef\grothLast{#2}}
\newcommand{\foldertitle}[1]{\gdef\grothFoldertitle{#1}}
\gdef\grothFolder{?}\gdef\grothBatch{?}\gdef\grothFirst{?}\gdef\grothLast{?}\gdef\grothFoldertitle{}

% --- Le résumé --------------------------------------------------------------
%
% Il ouvre la lecture modernisée, sous le titre « Résumé », et il est composé
% en petit. Ce n'est pas un ornement : c'est le seul passage du document qui
% ne soit pas des mathématiques, et un lecteur doit pouvoir le reconnaître
% avant de l'avoir lu — puis le sauter s'il connaît déjà le sujet, ou s'y
% arrêter s'il ne le connaît pas. Le filet qui le suit marque la reprise du
% corps.

\newenvironment{resume}
  {\small\setlength{\parskip}{0.45em}}
  {\par\medskip\hrule\medskip}

% --- Mots-clés ---------------------------------------------------------------
%
% En anglais, en clôture du résumé de la lecture modernisée : le vocabulaire
% moderne sous lequel chercher ce que les feuillets construisent. Le manifeste
% du site (`npm run manifest`) les extrait pour étiqueter la cote entière —
% cette ligne est la source unique des tags, il n'y a pas de fichier à côté.
\newcommand{\keywords}[1]{\par\smallskip\noindent
  {\footnotesize\textsc{Keywords} --- \textit{#1}}\par}

% --- L'appareil critique ----------------------------------------------------

% Le numéro de feuillet, en marge. C'est l'ancre qui permet de régler un
% désaccord sur une formule en regardant *un* feuillet plutôt qu'une chemise
% de six cents. Le site s'en sert aussi pour tourner les pages du fac-similé
% quand on fait défiler la transcription.
\newcommand{\leaf}[1]{%
  \marginnote{\footnotesize\color{gray!70}\textsf{#1}}%
  \label{leaf:#1}\ignorespaces}

% Illisible : on ne devine pas. Un mot inventé qui se lit comme les autres est
% le pire résultat possible.
\newcommand{\ill}{\textcolor{red!60!black}{[\,\ldots\,]}}

% Lecture douteuse : le mot est donné, et signalé comme incertain.
\newcommand{\uncertain}[1]{\uline{#1}}

% Ajout de l'éditeur — une lettre manquante, un mot sous-entendu.
\newcommand{\add}[1]{\textcolor{blue!50!black}{[#1]}}

% Biffé par Grothendieck lui-même. Ce qu'il a rayé fait partie du document :
% une rature dit souvent où la pensée a changé de direction.
\newcommand{\struck}[1]{\sout{#1}}

% Note du transcripteur, sur l'état matériel du feuillet.
\newcommand{\note}[1]{\par\smallskip{\small\color{gray!60!black}\textit{[#1]}}\par\smallskip}

% Note marginale de Grothendieck, distincte du corps du texte.
\newcommand{\marginal}[1]{\par\smallskip\noindent\hspace{1em}%
  {\small\color{gray!60!black}#1}\par\smallskip}

% --- Titre ------------------------------------------------------------------

\AtBeginDocument{%
  \begin{center}
    {\large\bfseries Fonds Alexandre Grothendieck --- cote n\textsuperscript{o}~\grothFolder}\\[2pt]
    {\itshape\grothFoldertitle}\\[4pt]
    {\small feuillets \grothFirst--\grothLast\ (lot \grothBatch)}\\[2pt]
    {\footnotesize\color{gray!60!black}Transcription établie d'après le fac-similé numérisé
     par l'Université de Montpellier.\\
     Les lectures incertaines sont soulignées, les illisibles notées [\,\ldots\,],
     les ajouts entre crochets.}
  \end{center}
  \vspace{1em}}

\endinput


\folder{66}
\batch{1}
\pages{1}{4}
\dating{1964}
\watermark{Édition de démonstration}
\foldertitle{Problèmes : notes manuscrites (1964)}

\begin{document}

\page{1}

\emph{En haut à gauche, encadré :} 1964.

\emph{En haut à droite :} Numéros cerclés sont sans doute difficiles.

\section*{Travaux de thèses}

\note{sous le titre, un sous-titre a été rayé au point d'être illisible : \ill}

(NB Les sujets 1 à 10 et 20, 21, \add{26, 27, 28} sont de nature essentiellement
cohomologique. 10, 19, 23, 24, 25 relèvent de la théorie des intersections ;
11, 12, 13 se rattachent à la th.\ de Néron. 14, 15, 22 th.\ des topos\ldots\
17, 18 yoga de représentabilité, formalisme général \ill)

\marginal{$*$ [Relation avec cohomologie cristalline]}

\note{les numéros que Grothendieck a entourés sont donnés ici entre
parenthèses, sous la forme (7) ; ce sont ceux qu'annonce la ligne de tête}

\begin{enumerate}
\item[1)] Topologie et cohomologies \ill\ fppf : formalisme général \ill

\item[(2)] \emph{(la ligne est laissée vide, un simple trait horizontal la
tient)} : dualité (\underline{locale}, globale).

\marginal{[subordonné à 1), 11)]}

\item[(3)] Partie $p$-primaire dans formule de Ogg--Chafarevitch

\marginal{[subordonné à 2)) donc à 1) et 11)]}

\item[(4)] Formule de RR en cohomologie étale, coeff.\ discrets. Relation avec
représentations d'Artin supérieures. [Formule des $\chi$ pour revêtement étale
d'un schéma propre\ldots]

\item[5)] Formule des $\chi$ pour V.A.\ avec anneau d'opérateurs (variante de
Ogg--Chafarevitch).

\item[6)] Formules de Lefschetz explicites (coeff.\ discrets étales, ou
« continus » cohérents). Représentations d'Artin du point de vue cohérent\ldots

\item[(7)] Cohomologie des schémas de type fini sur $\mathbb{Z}$ ; théorèmes de
dualité globale satisfaisants (améliorant Artin).

\item[(8)] Formalisme de dualité, du point de vue \struck{complexes}
\add{opérateurs} différentiels. Coefficients à connexion intégrable\ldots\ Lien
avec « cristaux ».

\item[(9)] Dualité analytique\ldots

\item[(10)] Riemann--Roch et \struck{pb} problèmes connexes (faisceaux de
Modules), analytiques ou algébriques. (Cf.\ \uncertain{SGA 6})

\item[(11)] Structures et groupes de la catégorie des schémas normaux.
Foncteur de Néron relatif à un anneau de val.\ discrète.
\end{enumerate}

\page{2}

\begin{enumerate}
\item[12)] Mise au point, et extensions diverses, de la théorie des modèles de
Néron.

\marginal{commencé par Raynaud}

\item[(13)] Théorèmes asymptotiques pour les foncteurs de Greenberg, relatifs à
un anneau local noethérien complet ou hensélien.

\marginal{\ill\ depuis les résultats d'approximation d'Artin}

\item[14)] Structures des Topos (questions d'existence de foncteurs
fibres\ldots). Structure des champs. champs représentables\ldots

\item[(15)] Théories du \struck{type} d'homotopie des Topos, incluant une
définition des $\pi_{i}$ à \struck{coefficients} morceaux infinitésimaux.
Th.\ d'Hurewicz. \struck{Suites} exactes d'homotopie. Th.\ de comparaison,
\add{« livrai »} \struck{\ill}

\marginal{fait partiellement par \uncertain{M.\ Artin}. Réflexions de Quillen
--- en train}

\item[(16)] Théorie de spécialisation du \add{gr.} gr.\ p.\ fondamental, pour
une famille de courbes.

\marginal{contre-exemple de Artin, schéma \ill}

\item[17)] Séparation des \struck{foncteurs} : classes de correspondance\ldots\
Cas de représentabilité et de non-représentabilité du
$\underline{\mathrm{Pic}}_{X/S}$, $X$ de dim.\ \struck{\ill} relative 1 sur $S$.

\marginal{résolu par Raynaud}

\item[(18)] Question des passages aux quotients dans cas arbitraires et de
l'effectivité de données de descente plates.
\struck{Construction générale de corps de \ill}

\item[(19)] Théorie d'amplitude pour cycles de dim.\ quelconque.
\end{enumerate}

\page{3}

\begin{enumerate}
\item[(20)] Vanishing theorems. Cas des espaces homogènes.

\item[(21)] $H^{*}(X, F) \to H^{*}(X_{0}, F_{0})$ ($X_{0}$ fermé dans $X$
\add{coh.\ étale}, quand est-ce un isom ? Suffit-il que ce soit vrai en dim 0 ?
Cas $X = \mathrm{Spec}\,A$, $X_{0} = \mathrm{Spec}\,A/J$, $A$ $J$-adiquement
séparé complet. Cas de la localisation étale le long de $X_{0} \subset X$, $X$
affine.

\item[22)] Structure des catégories abéliennes avec produit tensoriel $\pm$
rigide, en termes de représentations linéaires de groupes
\struck{proalgébriques}, et de représentations de gerbes\ldots\ (préliminaire
algébrique à la théorie des motifs).

\marginal{travail Saavedra en train}
\end{enumerate}

\note{c'est le programme tannakien : la thèse de Saavedra Rivano,
\emph{Catégories tannakiennes}, est ici notée en cours}

\emph{Encadré :} Pour mémoire

\begin{itemize}
\item Résolution des singularités.
\item Pureté cohomologique.
\end{itemize}

\emph{Une accolade réunit ces deux lignes et deux flèches en partent :}
$\Rightarrow$ dualité locale $=$ bidualité, pour le top.\ étale ;
$\Rightarrow$ th.\ de finitude pour morphismes \add{de type fini} etc.\ etc.

Changement de base régulier \add{en cohomologie}.

Dimension cohomologique des morphismes affines.

\begin{itemize}
\item Théorèmes de finitude genre Grauert.
\item Théories \struck{\ill} type rigide-analytiques.
\end{itemize}

\section*{Questions dans SGA 1962}

\begin{enumerate}
\item[23)] Théorie de Chow pour schémas \struck{réguliers} ayant un faisceau
inversible ample (équivalence rationnelle des cycles). L'image inverse par
$f : X \to Y$ doit être définie, sous seule condition que $Y$ soit régulier.

\item[(24)] 1\textsuperscript{o} Sur une variété projective lisse connexe de
dim \ill\ sur $k$ corps algébriquement clos, un 1-cycle num.\ équiv.\ à 0
est-il \ill\ équivalent à zéro ? En dim \ill, quand \ill\ équiv.\ num.\
$\approx$ équiv.\ hom. ?

\marginal{Réponse négative Griffiths pour \ill}
\end{enumerate}

\note{les deux dernières lignes de la page courent jusqu'au bord inférieur du
cliché et n'y sont plus qu'en partie lisibles ; ce qui manque est signalé et
non restitué}

\page{4}

\begin{enumerate}
\item[(25)] Le produit de deux cycles alg.\ équiv.\ à 0 est-il rat.\ équiv.\ à
0 ? (« Pb des \uncertain{courbes} »). Un 0-cycle qui est alg.\ équiv.\ à 0
est-il rationnellement équivalent à 0 ? Question de l'équivalence rationnelle
des 0-cycles sur une surface.

\marginal{faux, cf Mumford\ldots}

\item[26] Théorie de Rosenlicht relative (sur base quelconque).

\item[27] Adèles des schémas de dim.\ quelconque. Jacobiennes locales.

\item[(28)] Construction d'une théorie abstraite des motifs sur préschémas de
type fini sur $\mathbb{Z}$ (construction catégorique générale).
\end{enumerate}

\note{la liste s'arrête à 28 et le reste de la page est vierge}

\end{document}
